\documentclass[11pt,a4paper]{article}

% -------------------------------------------------------------------
% LinkedIn article draft (separate from docs/article.tex)
% Compile: pdflatex + bibtex + pdflatex + pdflatex
% -------------------------------------------------------------------

\usepackage[T1]{fontenc}
\usepackage[utf8]{inputenc}
\usepackage{palatino}
\usepackage{microtype}
\usepackage{geometry}
\geometry{
  top=2.4cm,
  bottom=2.4cm,
  left=2.6cm,
  right=2.6cm,
  headheight=14pt
}

\usepackage[dvipsnames,table]{xcolor}
\definecolor{brandblue}{HTML}{0A66C2}
\definecolor{textdark}{HTML}{1F2937}
\definecolor{textmuted}{HTML}{6B7280}
\definecolor{okgreen}{HTML}{2E7D32}
\definecolor{warnorange}{HTML}{D97706}
\definecolor{softbg}{HTML}{F5F7FA}

\usepackage{titlesec}
\titleformat{\section}
  {\large\bfseries\color{brandblue}}{}{0em}{}
  [\vspace{-0.3em}\rule{\linewidth}{1.2pt}\vspace{0.15em}]
\titleformat{\subsection}
  {\normalsize\bfseries\color{textdark}}{}{0em}{}

\usepackage{parskip}
\setlength{\parskip}{0.5em}
\setlength{\parindent}{0pt}

\usepackage{enumitem}
\setlist[itemize]{leftmargin=1.4em,itemsep=0.2em,topsep=0.2em}
\setlist[enumerate]{leftmargin=1.6em,itemsep=0.2em,topsep=0.2em}

\usepackage{booktabs}
\usepackage{tabularx}
\usepackage{array}
\renewcommand{\arraystretch}{1.2}

\usepackage{tcolorbox}
\tcbuselibrary{skins,breakable}
\newtcolorbox{insightbox}[1][]{
  enhanced,
  breakable,
  colback=brandblue!5,
  colframe=brandblue,
  fonttitle=\bfseries\small,
  title={#1},
  left=6pt,right=6pt,top=4pt,bottom=4pt,arc=3pt
}
\newtcolorbox{scopebox}[1][]{
  enhanced,
  breakable,
  colback=warnorange!8,
  colframe=warnorange,
  fonttitle=\bfseries\small,
  title={#1},
  left=6pt,right=6pt,top=4pt,bottom=4pt,arc=3pt
}

\usepackage{amsmath}
\usepackage[
  colorlinks=true,
  linkcolor=brandblue,
  citecolor=brandblue,
  urlcolor=brandblue
]{hyperref}

\usepackage{fancyhdr}
\pagestyle{fancy}
\fancyhf{}
\fancyhead[L]{\small\color{textmuted}\textit{ThingsNode -- S.\,M.\,Pasindu Upeksha Senevirathne}}
\fancyhead[R]{\small\color{textmuted}March 2026}
\fancyfoot[C]{\small\color{textmuted}\thepage}
\renewcommand{\headrulewidth}{0.3pt}

\usepackage[
  backend=bibtex,
  style=numeric,
  sorting=none,
  maxbibnames=3
]{biblatex}
\addbibresource{article.bib}

\newcommand{\artifact}[1]{\texttt{#1}}
\newcommand{\metric}[1]{\textbf{#1}}
\newcommand{\todoitem}[1]{\textcolor{warnorange}{#1}}

\begin{document}

\begin{center}
  {\fontsize{21}{27}\selectfont\bfseries\color{textdark}
  From Initial Plan to Production Model:\\
  Physics-Based PV Expected Generation for IoT SCADA}\\[8pt]
  {\large\color{brandblue}
  A practical engineering journey from project plan to production deployment,\\
  based on repository code, configs, data, and output artifacts}\\[10pt]
  {\small\color{textmuted}
  S.\,M.\,Pasindu Upeksha Senevirathne \quad|\quad ThingsNode \quad|\quad March 2026 \quad|\quad LinkedIn}
\end{center}

\begin{insightbox}[Scope of this article]
This document is intentionally evidence-first. It is written from repository artifacts:
\artifact{docs/}, \artifact{scripts/}, \artifact{config/plant\_config.json},
\artifact{data/}, and \artifact{output/}. The final reference model is
\artifact{scripts/python\_physics/physics\_model.py}; other irradiance-to-generation paths are documented as trial-and-error work.
This project was developed individually during my internship at ThingsNode.
\end{insightbox}

% ===================================================================
\section{Structure and Conventions}
% ===================================================================

This article follows a practical engineering narrative with artifact-backed claims.

\textbf{Document structure used in Overleaf}
\begin{enumerate}
  \item Problem and business need
  \item What was built (production path)
  \item What was intentionally out of scope
  \item What was tested and discarded
  \item Evidence from generated outputs
  \item Market comparison and positioning
  \item Research alignment and technical rationale
  \item Product strengths, weaknesses, and fit
  \item Adoption playbook and roadmap
\end{enumerate}

\textbf{Conventions used}
\begin{itemize}
  \item \textbf{Evidence-tagged claims:} metrics are tied to explicit files (for example, \artifact{output/js\_validation\_results.json}).
  \item \textbf{Model path separation:} reference model vs helper/experimental paths are never mixed.
  \item \textbf{Business relevance first:} each technical choice is explained with product impact (accuracy, speed, cost, operability).
  \item \textbf{Caveat-first reporting:} when a metric has window or pipeline bias, that caveat is stated beside the metric.
  \item \textbf{Currency convention:} when monetary values are referenced, USD is used for international readability.
  \item \textbf{Client anonymization:} plant owner name is intentionally omitted in this public version.
\end{itemize}

% ===================================================================
\section{Problem and Why It Matters}
% ===================================================================

Solar operations teams need expected generation at telemetry cadence, not only annual planning reports.
Without that baseline, the following are hard to detect in time: soiling drift, inverter underperformance,
and weather-normalized PR deterioration.

The project started from \artifact{docs/project\_plan.md}, then evolved through multiple implementation paths.
The key product question became:

\begin{quote}
\textit{Can we produce an expected-generation engine that is physics-consistent, auditable, and operationally deployable in an IoT pipeline?}
\end{quote}

\textbf{Project deployment context and data context}
\begin{itemize}
  \item Primary dataset context: Dewanappiyatissa School PV site, Sri Lanka
  \item Site coordinates: latitude 8.342368984714714, longitude 80.37623529556957
  \item Repository configuration context: 55\,kW AC rooftop system (\artifact{config/plant\_config.json})
  \item Geometry in config: 8 orientations, 216 modules, 552.096\,m$^2$ total module area
\end{itemize}

% ===================================================================
\section{What We Built (Production Reference Path)}
% ===================================================================

\subsection{Reference Model: \artifact{physics\_model.py}}

The retained model is \artifact{scripts/python\_physics/physics\_model.py}, built with pvlib \cite{holmgren2018pvlib}.
It supports hourly and daily irradiance inputs (daily can run through synthetic-hourly reconstruction, or a simpler daily branch).

\textbf{Core chain in execution order}
\begin{enumerate}
  \item Source ingestion: Solcast API, NASA POWER hourly/daily, CSV hourly/daily
  \item Solar position via pvlib location/SPA \cite{reda2004solar}
  \item POA transposition with Perez model \cite{perez1990set}
  \item Far-shading multiplier
  \item AOI + IAM interpolation table
  \item SAPM cell temperature model \cite{king2004photovoltaic}
  \item DC conversion with temperature correction
  \item Sequential DC loss chain from configured loss terms
  \item DC-to-AC conversion (curve-based or flat, controlled by config)
  \item Plant-level clipping, then AC wiring loss
\end{enumerate}

\textbf{Reference formulation}
\[
P_{AC}(t)=
\left(
\min\left[
\sum_i P_{DC,i}(t)\cdot \eta_{inv}(P_{DC,i}(t)),\;
P_{AC,\max}
\right]
\right)\cdot(1-L_{AC})
\]
where each orientation $i$ uses Perez POA, IAM, SAPM thermal correction, and configured DC losses.

\subsection{Operational Orchestration: \artifact{daily\_predictor.py}}

\artifact{scripts/python\_physics/daily\_predictor.py} runs daily and merges weather fields through source fallback chains.
It writes append-mode hourly/daily/monthly outputs under:
\artifact{output/hourly/}, \artifact{output/daily/}, \artifact{output/monthly/}.

\begin{insightbox}[Important implementation nuance]
\artifact{daily\_predictor.py} uses a simplified inline physics variant (flat inverter efficiency for speed).
The reference math remains \artifact{physics\_model.py}. This separation is explicit in current docs
(\artifact{docs/model\_methods.md}, \artifact{docs/article\_evidence\_pack.md}).
\end{insightbox}

\subsection{Chosen Production Deployment: AWS IoT Core + Lambda}

The final production choice is the Python reference model on AWS, not a JavaScript edge runtime:
\begin{enumerate}
  \item telemetry enters AWS IoT Core,
  \item IoT rule triggers AWS Lambda,
  \item Lambda executes \artifact{scripts/python\_physics/physics\_model.py},
  \item expected generation is published back to IoT topics and dashboards.
\end{enumerate}

This decision prioritized accuracy over speed. JavaScript paths were explored and benchmarked,
but were not selected as the production baseline.

\subsection{Configuration and Auditability}

The design is config-driven, not hardcoded:
\begin{itemize}
  \item geometry and orientation count,
  \item module and inverter parameters,
  \item loss factors and IAM table,
  \item timezone and defaults for weather fallbacks.
\end{itemize}

This enables reproducibility per plant without changing core code.

% ===================================================================
\section{What We Intentionally Did Not Build (Yet)}
% ===================================================================

\begin{scopebox}[Deliberate scope boundaries]
These are known boundaries, not accidental omissions. Each was assessed against
delivery timeline, required data, and operational payoff.
\end{scopebox}

\begin{itemize}
  \item \textbf{No near-shading geometry engine}: no row-to-row or 3D obstruction simulation in this codebase.
  \item \textbf{No spectral correction module}: no explicit spectral mismatch correction layer.
  \item \textbf{No annual degradation tracking}: LID exists, long-term degradation trend module is not implemented.
  \item \textbf{No battery/self-consumption dispatch model}: expected gross generation is modeled, not full site energy economics.
  \item \textbf{No automatic self-calibration of losses}: losses are static from config/PVsyst extraction process.
  \item \textbf{No completed annual PVsyst parity study}: still pending as a formal bankability-grade validation.
\end{itemize}

These limits are documented in \artifact{docs/model\_methods.md} and \artifact{docs/article\_research.md}.

% ===================================================================
\section{Trial-and-Error Work and Why It Was Discarded}
% ===================================================================

\subsection{Simplified Daily Model Branch}
\begin{itemize}
  \item File path: \artifact{calculate\_daily\_energy\_simple()} in \artifact{physics\_model.py}
  \item Benefit: very fast daily estimate
  \item Rejection reason: weak hourly fidelity, hardcoded POA factor behavior, reduced physical realism
\end{itemize}

\subsection{Synthetic Hourly Reconstruction from Daily Totals}
\begin{itemize}
  \item File path: \artifact{generate\_synthetic\_hourly\_from\_daily()}
  \item Benefit: allows full hourly chain when only daily totals are available
  \item Limitation: clear-sky-shaped intraday profile can diverge heavily on convective cloud days
\end{itemize}

\subsection{JS Surrogate Regression}
\begin{itemize}
  \item Files: \artifact{scripts/js\_surrogate/fit\_surrogate.py}, \artifact{surrogate\_model.js}
  \item Equation: $a_0 + a_1GHI + a_2GHI^2 + a_3DNI + a_4(GHI\Delta T)+a_5(GHIv)$
  \item Outcome: very fast, but poor hourly behavioral fidelity
  \item Evidence: on held-out test (\artifact{output/test\_evaluation\_metrics.json}), \metric{$R^2=0.420$}, MAE 9.44\,kW, MAPE 129.7\%
\end{itemize}

\subsection{Standalone DNI Decomposition Script}
\begin{itemize}
  \item File: \artifact{scripts/predict\_dni.py}
  \item Value: diagnostic consistency check
  \item Status: not required for final expected-generation pipeline, archived as utility
\end{itemize}

% ===================================================================
\section{Evidence from Output Artifacts}
% ===================================================================

\subsection{Internal Parity: JS Physics Approximation vs Python Labels (Exploratory)}

These internal parity metrics come from the repository's 55\,kW configuration artifacts.

\begin{table}[h]
\centering
\small
\caption{Calibration and validation of JS physics approximation}
\begin{tabularx}{\linewidth}{>{\raggedright\arraybackslash}p{4.2cm} X}
\toprule
Metric source & Result \\
\midrule
\artifact{output/js\_physics\_calibration.json} &
Baseline RMSE 2.2449\,kW \textrightarrow{} Final RMSE 1.4340\,kW; MAE 0.6723\,kW; $R^2=0.99233$; energy error 1.491\%; 112 iterations, 13,677 evaluations \\
\artifact{output/js\_validation\_results.json} &
2,500,000 samples; RMSE 1.4577\,kW (2.65\% of 55\,kW), MAE 0.6763\,kW, $R^2=0.99214$, MAPE 7.76\%, energy error 1.456\% \\
\bottomrule
\end{tabularx}
\end{table}

\subsection{Held-Out 2024 Test Set (JS Models vs Python-Labeled Test Data)}

\begin{table}[h]
\centering
\small
\caption{Held-out test metrics from \artifact{output/test\_evaluation\_metrics.json}}
\begin{tabular}{>{\raggedright\arraybackslash}p{3.0cm} >{\raggedleft\arraybackslash}p{1.5cm} >{\raggedleft\arraybackslash}p{1.8cm} >{\raggedleft\arraybackslash}p{1.8cm} >{\raggedleft\arraybackslash}p{2.0cm}}
\toprule
Model & $R^2$ & MAE (kW) & MAPE & Energy error \\
\midrule
JS Physics & 0.9685 & 1.3387 & 20.09\% & +5.47\% \\
JS Surrogate & 0.4200 & 9.4375 & 129.70\% & +1.50\% \\
\bottomrule
\end{tabularx}
\end{table}

\textbf{Provenance caveat:} these metrics are produced by \artifact{scripts/shared/evaluate\_models.py},
which uses an inline JS-physics evaluator path for test execution.

\subsection{Against Real Meteocontrol Daily Actuals (December 2025)}

From \artifact{output/benchmark\_results.csv}, only two days contain overlapping non-null actuals:
\begin{itemize}
  \item 11 Dec 2025: Actual 345.36\,kWh vs Python 224.53\,kWh
  \item 12 Dec 2025: Actual 75.982\,kWh vs Python 263.41\,kWh
\end{itemize}

Overlap-only totals (2 days):
\begin{itemize}
  \item Python: +15.81\%
  \item JS Physics: +27.03\%
  \item JS Surrogate: +18.29\%
\end{itemize}

\textbf{Critical reporting caveat:}
\artifact{output/benchmark\_metrics.json} reports 5-day totals with inflated total-error percentages because
predictions exist for days lacking non-null actuals in the merged comparison window.

\subsection{Irradiance Source Comparison on 10-day Matched Window}

From \artifact{output/accuracy\_report.txt}:
\begin{itemize}
  \item NASA POWER: MAE 96.2\,kWh/day, bias +32.7\,kWh/day, $R^2=-0.68$
  \item Solcast: MAE 97.4\,kWh/day, bias +33.9\,kWh/day, $R^2=-0.71$
  \item Open-Meteo: MAE 111.4\,kWh/day, bias +47.9\,kWh/day, $R^2=-1.55$
  \item ERA5: MAE 96.8\,kWh/day, bias +33.3\,kWh/day, $R^2=-0.70$
\end{itemize}

This window demonstrates input-data limitation during monsoon-like cloud conditions:
all sources bias high, so all models over-predict.

% ===================================================================
\section{What Market Solutions Have and Do Not Have}
% ===================================================================

\subsection{Where the market is strong}
\begin{itemize}
  \item \textbf{Inverter-vendor portals (SolarEdge/SMA/Huawei class):}
  strong hardware-native diagnostics and device-level alarms.
  \item \textbf{PVsyst:} bankable yield planning standard, mature loss accounting workflow.
  \item \textbf{Irradiance specialists (for example Solcast):}
  high-quality weather and irradiance products as input layers.
\end{itemize}

\subsection{Where gaps remain for mixed-vendor IoT SCADA fleets}
\begin{itemize}
  \item Vendor lock-in across data logger, inverter, and analytics stacks
  \item Limited transparency of ``expected generation'' internals in black-box dashboards
  \item Limited transparent expected-generation tooling that is portable across mixed-vendor fleets
  \item Cost friction for small and mid-size operators that need auditable expected-generation baselines
\end{itemize}

\begin{table}[h]
\centering
\small
\caption{Capability comparison (product view)}
\begin{tabularx}{\linewidth}{>{\bfseries\raggedright\arraybackslash}p{3.5cm} >{\raggedright\arraybackslash}p{3.2cm} >{\raggedright\arraybackslash}p{3.2cm} >{\raggedright\arraybackslash}X}
\toprule
Capability & Our reference path & Typical vendor portal & Open-source toolkit only \\
\midrule
Transparent model internals & Full (script + config) & Usually limited & Full \\
Mixed-vendor portability & High & Low-to-medium & High \\
Real-time serverless fit (AWS IoT Core + Lambda) & High & Vendor specific & Medium (integration effort) \\
Bankability workflow & Partial / pending annual parity & Medium & Depends on implementation \\
Cost control and self-hosting & High & Medium-to-low & High \\
\bottomrule
\end{tabularx}
\end{table}

% ===================================================================
\section{How This Aligns with Existing Research and Tools}
% ===================================================================

This implementation direction is consistent with known literature and tools:
\begin{itemize}
  \item pvlib as a practical open modeling foundation \cite{holmgren2018pvlib}
  \item Perez transposition preference for non-trivial orientation geometry \cite{perez1990set,gueymard2009direct}
  \item SAPM thermal modeling benefits for irradiance+wind-driven cell temperature \cite{king2004photovoltaic}
  \item Limits of pure data-driven short-horizon models without strong site-specific history \cite{diagne2013review,mellit2010artificial}
  \item Transfer from physics ``oracle'' to faster approximation using large synthetic datasets \cite{inman2013solar,tobin2017domain}
\end{itemize}

Optimization choice (differential evolution) also matches the non-convex, interacting parameter landscape used in calibration \cite{storn1997differential}.

% ===================================================================
\section{Strengths, Weaknesses, and Product Position}
% ===================================================================

\subsection{Strengths}
\begin{itemize}
  \item Physics-first, auditable expected-generation engine
  \item Strong geometry handling for multi-orientation rooftops
  \item Config-driven reproducibility and operational script ecosystem
  \item Demonstrated fast approximation path when edge constraints demand it
  \item Clear artifact trail for QA and documentation
\end{itemize}

\subsection{Weaknesses and risks}
\begin{itemize}
  \item No near-shading simulation layer
  \item No completed annual PVsyst parity report in-repo
  \item Benchmark windows against actuals are currently short and weather-sensitive
  \item Helper-script model variants can confuse reporting if provenance is not explicit
\end{itemize}

\subsection{Why this can be better than a generic product}
\begin{itemize}
  \item Full inspectability of assumptions and losses
  \item No forced hardware ecosystem coupling
  \item Ability to tune to each plant from first principles and explicit config
\end{itemize}

\subsection{Where this can be worse than a generic product}
\begin{itemize}
  \item Requires engineering ownership; no one-click ``managed analytics''
  \item Lacks vendor-native deep hardware diagnostics
  \item Accuracy remains bounded by weather-input fidelity during severe local cloud events
\end{itemize}

% ===================================================================
\section{Who This Is For and How to Use It}
% ===================================================================

\textbf{Best-fit users}
\begin{itemize}
  \item Solar O\&M teams needing transparent expected-generation baselines
  \item IoT/SCADA engineering teams integrating expected generation into telemetry pipelines
  \item EPC and analytics teams that already extract parameters from PVsyst-style commissioning outputs
\end{itemize}

\textbf{Not the right fit}
\begin{itemize}
  \item Pure financial bankability workflows needing certified end-to-end vendor tooling
  \item Teams requiring string-level hardware diagnostics from inverter-native data channels
\end{itemize}

\textbf{Adoption playbook}
\begin{enumerate}
  \item Populate \artifact{config/plant\_config.json} for a target site.
  \item Run reference physics from \artifact{scripts/python\_physics/physics\_model.py}.
  \item Generate operational predictions via \artifact{daily\_predictor.py}.
  \item Compare with actuals using \artifact{scripts/shared/compare\_accuracy.py}.
  \item Deploy the Python reference path in AWS Lambda with AWS IoT Core rule triggering.
\end{enumerate}

% ===================================================================
\section{Business Need and ROI Logic}
% ===================================================================

The product value is less about ``perfect forecasting'' and more about \textit{operational decision speed}.
When expected generation is transparent and continuously available:
\begin{itemize}
  \item underperformance can be detected earlier,
  \item cleaning and maintenance timing becomes evidence-based,
  \item PR discussions become weather-normalized instead of nameplate-normalized,
  \item performance conversations become auditable.
\end{itemize}

For a 55\,kW site, even small persistent bias (for example 3--5\% unmanaged loss) can accumulate into meaningful monthly energy and revenue leakage.

In this project, one consistent operational observation was that model behavior improves substantially when irradiance input quality is good; many sites perform as expected when the weather inputs are accurate and well-aligned to local conditions.

% ===================================================================
\section{Roadmap}
% ===================================================================

\begin{itemize}
  \item Complete annual PVsyst parity validation for the deployed plant(s)
  \item Add automated overlap-aware benchmarking logic (avoid inflated total-error windows)
  \item Add near-shading and long-term degradation modules
  \item Implement self-calibration strategy for slowly drifting loss factors
  \item Publish a hardened deployment profile separating reference and helper model paths
\end{itemize}

% ===================================================================
\section{Conclusion}
% ===================================================================

This project now has a clear technical and product narrative:
\begin{itemize}
  \item one retained reference model (\artifact{physics\_model.py}),
  \item multiple documented trial paths with explicit rejection reasons,
  \item measurable internal validation artifacts,
  \item and an honest boundary on real-world accuracy when irradiance inputs are wrong.
\end{itemize}

That combination -- physics rigor, operational pragmatism, and artifact-level transparency -- is the core business differentiator.

\vspace{0.6em}
\noindent\rule{\linewidth}{0.4pt}
{\small\color{textmuted}
Primary evidence files used in this article include:
\artifact{docs/project\_plan.md},
\artifact{docs/model\_methods.md},
\artifact{docs/article\_evidence\_pack.md},
\artifact{docs/article\_research.md},
\artifact{scripts/python\_physics/physics\_model.py},
\artifact{scripts/shared/evaluate\_models.py},
\artifact{output/js\_physics\_calibration.json},
\artifact{output/js\_validation\_results.json},
\artifact{output/test\_evaluation\_metrics.json},
\artifact{output/benchmark\_metrics.json},
\artifact{output/benchmark\_results.csv},
\artifact{output/accuracy\_report.txt}.}

\newpage
\printbibliography[title={References}]

\end{document}
